\begin{abstract}
Paper~B2 reports an expanded empirical investigation executed under the Invariance-Based Architectural Investigation Protocol v1.0 (``Protocol v1.0''), defined independently in Paper~0 as the scientific instrument. Paper~B1 reported a three-system pilot execution; Paper~B2 applies the same frozen protocol to a governance/authorization cohort.

The objective is cohort-bounded: to produce the protocol-required artifacts (Baseline Observation Records, execution surface matrix, derived object registry, measurement registry, comparative dataset, retained classifications, and canonical cohort conclusion) and to evaluate---within the evaluated cohort only---whether a candidate Reasoning\,$\rightarrow$\,Legitimacy\,$\rightarrow$\,Execution (RLE) boundary is observed under the protocol's operationalization. This paper does not modify Protocol v1.0, does not rewrite Paper~B1, and avoids universal invariant claims.

\textbf{Cohort (candidate systems).} OPA+Gatekeeper; Kubernetes RBAC/Admission; AWS IAM; HashiCorp Vault; Envoy \texttt{ext\_authz}; Istio \texttt{AuthorizationPolicy}; OpenFGA; Cedar / Amazon Verified Permissions; Google Zanzibar.

\textbf{Results (cohort-bounded).} The canonical cohort conclusion is recorded in the machine-readable cohort-conclusion artifact generated from retained classifications and the frozen I5 decision rule.

\textbf{Artifacts.} Machine-readable B2 artifacts are stored under \texttt{investigations/b2-governance-cohort/}; the canonical cohort conclusion is \texttt{results/b2-governance-cohort-i5.cohort-conclusion.json}.
\end{abstract}
